\section{Introduction}

A smooth connected complex projective variety $X$ is \emph{Fano} if
its anticanonical line bundle $\omega_X^{-1}$ is ample. A vector bundle
$E$ is \emph{nef} if the tautological line bundle
$\mathcal O_{\mathbb P(E)}(1)$ has nonnegative degree on every curve.
We use the quotient convention
$\mathbb P(E)=\operatorname{Proj}\operatorname{Sym}E$.
A variety is \emph{rational homogeneous} if it is isomorphic to $G/P$,
where $G$ is a connected semisimple complex algebraic group and $P$ is
a parabolic subgroup.

The tangent bundle of a rational homogeneous variety is generated by
global sections and is therefore nef. Mori's characterization of
projective space by ampleness of its tangent bundle
\cite{Mori1979} makes the corresponding question for nef tangent
bundles natural. The Campana--Peternell conjecture asks whether every
smooth complex Fano variety with nef tangent bundle is rational
homogeneous \cite{CampanaPeternell1991}. Thus it predicts that a
numerical positivity condition determines a variety with a transitive
algebraic group action. We prove the conjecture in dimension six.

\begin{theorem}\label{thm:main}
Every smooth connected complex projective Fano variety of dimension six
with nef tangent bundle is rational homogeneous.
\end{theorem}

The Picard number $\rho(X)$ is the rank of the group of divisor classes
modulo numerical equivalence, and the \emph{pseudoindex} is the least
anticanonical degree of a rational curve on $X$.
These two invariants organize much of the progress on the conjecture.
Campana and Peternell established it in dimensions at most three and
for fourfolds of Picard number greater than one
\cite{CampanaPeternell1991,CampanaPeternell1993}.
The remaining fourfold case was settled through Mok's and Hwang's
study of the tangent directions of minimal rational curves
\cite{Mok2002,Hwang2006}. Watanabe proved the fivefold case of Picard
number greater than one \cite[Theorem~1.2]{Watanabe2014}, and
Kanemitsu completed the fivefold classification
\cite[Theorem~0.2]{Kanemitsu2017}.

Kanemitsu also proved that a smooth complex Fano $n$-fold with nef
tangent bundle is rational homogeneous whenever $\rho(X)>n-5$
\cite[Theorem~0.2]{Kanemitsu2016}. For a sixfold, this covers every
$\rho(X)\geq2$, leaving only Picard number one. In that range,
standard characterization results reduce the problem to pseudoindices
four and five. Watanabe treated the pseudoindex-four sixfold case
\cite[Theorem~2.4]{Watanabe2021}. We include an independent intersection
argument for it; the principal new step is pseudoindex five, where a
characteristic-number constraint leads to the Grassmannian.

The difficulty is to pass from positivity of $T_X$ to a specific
projective model for the tangent directions of minimal rational curves.
Nefness gives strong smoothness and deformation properties, but these
properties do not by themselves identify the threefold of directions
that occurs in pseudoindex five. We extract additional information by
using the universal family in two ways: as a family of rational curves
over their parameter space, and as a family of threefolds over $X$.
The compatibility of these two maps yields a characteristic-number
identity strong enough to determine the projective variety of tangent
directions.

For a covering family of rational curves of least anticanonical degree,
the closure of the tangent directions to its members smooth at a general
point is called the \emph{variety of minimal rational tangents} (VMRT).
Characteristic classes and the geometry of minimal curves already
interact in Mok's work on one-dimensional varieties of minimal rational
tangents \cite{Mok2002}. Under an additional fourth-Betti-number
hypothesis, his argument uses the two rational-curve fibrations of the
universal family and characteristic-class inequalities for a rank-two
bundle on the parameter space to constrain the tangent image.
Hwang removed that hypothesis
\cite[Section~4]{Hwang2006}; the resulting classification is stated in
\cite[Theorem~15]{Mok2008}. A related structural approach studies
elementary Mori contractions, the contractions of extremal rays:
if all of these contractions of a Fano manifold are smooth
$\mathbb P^1$-fibrations, the manifold is a complete flag variety
\cite[Theorem~1.2]{OcchettaSolaCondeWatanabeWisniewski2017}.
These earlier results show how multiple curve fibrations can constrain
both characteristic classes and homogeneous structure.
The calculation here concerns a three-dimensional tangent image. It
uses the two fibrations to produce explicit linear intersection
identities, then combines them with the product identities furnished
by the projection formula.

\subsection{The argument}

Assume $\rho(X)=1$. Standard characterization theorems
leave pseudoindices four and five; the other possibilities are excluded
or give projective space $\mathbb P^6$ and the smooth quadric $Q^6$.
The proof is organized around the universal family of rational curves
of least degree. Its evaluation map and its projection to the parameter
space are both smooth. Along every fibre of the latter map, the
relative tangent bundle of evaluation is a sum of copies of
$\mathcal O_{\mathbb P^1}(-1)$.

For pseudoindex four, Kanemitsu's factorization theorem reduces the
problem to a variety admitting two fibrations. We rule out the possible
fibres by Hard Lefschetz and intersection theory. The last case is an
odd-dimensional quadric bundle: its orthogonal description forces a
vanishing characteristic number, whereas the second fibration would
make an integral divisor have nonintegral degree.

For pseudoindex five and a general point $x\in X$, the evaluation fibre
$F$ is a smooth threefold parametrizing minimal rational curves marked at
$x$. Sending a marked
curve to its tangent direction at $x$ gives a finite morphism
$F\to\mathbb P(T_x^*X)\simeq\mathbb P^5$. Its image is the variety of
minimal rational tangents, and the pullback $L$ of $\mathcal O(1)$ is
ample. At this point the degree of the map and the projective type of
its image are both undetermined. The central numerical step is
\[
 32-20u+3u^2=0,
 \qquad u=\frac{c_1(T_F)\cdot L^2}{L^3}.
\]
We derive this identity from the cohomology of the universal family,
Deligne's fixed-part theory \cite{Deligne1971}, and
Grothendieck--Riemann--Roch \cite{BorelSerre1958}. All resulting characteristic-number
constraints are written explicitly. A finite rational calculation then
extracts the quadratic relation.

This method separates two useful ingredients. The two fibrations give
linear relations among intersection numbers; integration along evaluation
also makes selected numbers factor as products. Multiplying these
relations forces a nonlinear constraint on the polarized fibre.
The proof explains both ingredients and gives an exact certificate for
the finite elimination.

Finally, this image spans $\mathbb P^5$. The sectional-genus
formula excludes $u=4$ and forces $L^3=3$ when $u=8/3$.
The variety of minimal rational tangents is consequently the Segre
threefold $\mathbb P^1\times\mathbb P^2\subset\mathbb P^5$.
The VMRT method of Hwang and Mok
\cite[Propositions~1.2.1--1.2.2]{HwangMok1999} supplies the
nondegeneracy criterion, and Mok's recognition theorem
\cite{Mok2008Recognition} identifies $X$ with
$\operatorname{Gr}(2,5)$, the Grassmannian of two-dimensional subspaces
of $\mathbb C^5$.

\subsection{Consequences of the classification}

The classification also has a consequence for compact K\"ahler manifolds
that need not be projective. A differential-geometric antecedent is
Mok's uniformization theorem for compact K\"ahler manifolds with
nonnegative holomorphic bisectional curvature \cite{Mok1988}.
For nef tangent bundles, the structure theorem of
Demailly--Peternell--Schneider relates the compact K\"ahler problem to
the Fano conjecture through the fibres of the Albanese map
\cite[Main Theorem and Theorem~3.14]{DemaillyPeternellSchneider1994}.
To state the consequence precisely, a holomorphic vector bundle
is \emph{nef in the analytic sense} when its tautological line bundle
admits smooth Hermitian metrics whose curvature is bounded below by
$-\varepsilon\omega$ for every $\varepsilon>0$, for a fixed Hermitian
form $\omega$ on its projectivization. This agrees with the preceding
numerical definition in the projective setting
\cite[Section~1.A and Definitions~1.2 and~1.9]{DemaillyPeternellSchneider1994}.
For a connected compact K\"ahler manifold $X$ with nef tangent bundle, put
\[
 \widetilde q(X)=
 \max\bigl\{h^1(X',\mathcal O_{X'}):
       X'\longrightarrow X\text{ is a connected finite \'{e}tale cover}\bigr\}.
\]
This maximum exists and is at most $\dim_{\mathbb C}X$
by the Demailly--Peternell--Schneider Albanese theorem
\cite[Proposition~3.9 and Theorem~3.14]{DemaillyPeternellSchneider1994}.

\begin{corollary}[Albanese fibres and uniformization]\label{cor:nef-tangent-uniformization}
Let $X$ be a connected compact K\"ahler manifold whose holomorphic
tangent bundle is nef in the analytic sense. Suppose
\[
 0\leq d=\dim_{\mathbb C}X-\widetilde q(X)\leq6.
\]
Then there is a connected finite \'{e}tale cover $X'\to X$ with the
following properties.
\begin{enumerate}
\item The Albanese map
      $\alpha:X'\to A=\operatorname{Alb}(X')$ is a smooth proper
      surjection with connected fibres, where $A$ is a complex torus of
      dimension $\widetilde q(X)$. Every fibre is rational homogeneous
      of dimension $d$, with a point allowed when $d=0$.
\item There is one such fibre $F$ for which $\alpha$ is a holomorphically
      locally trivial $F$-bundle. The ordinary universal cover of $X$
      is biholomorphic to
      \[
       \widetilde X\simeq F\times\mathbb C^{\widetilde q(X)}.
      \]
      It is therefore biholomorphic to a quasi-projective variety and
      is holomorphically convex. Its Remmert reduction is the projection
      \[
       F\times\mathbb C^{\widetilde q(X)}
       \longrightarrow\mathbb C^{\widetilde q(X)}.
      \]
\end{enumerate}
In particular, these conclusions hold whenever $\dim_{\mathbb C}X\leq6$.
\end{corollary}

The first conclusion applies the Fano classification fibrewise to the
Demailly--Peternell--Schneider fibration. The second also uses local
deformation rigidity of rational homogeneous manifolds, the
Fischer--Grauert theorem, and the universal-cover splitting theorem of
Claudon--H\"oring--Koll\'ar. The torus $A$ may be nonprojective, and the
finite-cover bundle may remain twisted: the product is asserted only
for the ordinary universal cover. We prove the corollary after the
sixfold theorem in Section~\ref{sec:conclusion}.

The Fano classification also gives a positive K\"ahler--Einstein metric
on every Fano manifold of positive dimension at most six with nef tangent bundle;
Corollary~\ref{cor:positive-ke} records the normalization of this metric.

\paragraph{Organization.}
Section~\ref{sec:minimal-curves} establishes the common geometry,
including the rational-curve fibration calculation in
Section~\ref{subsec:p1-cohomology}.
Section~\ref{sec:pseudoindex-four} settles pseudoindex four.
Sections~\ref{sec:characteristic-classes} and~\ref{sec:arithmetic}
derive the quadratic relation. Section~\ref{sec:grassmannian}
recognizes the Grassmannian, and Section~\ref{sec:conclusion}
combines the cases with Kanemitsu's Picard-number theorem and derives
Corollary~\ref{cor:nef-tangent-uniformization}.
Appendix~\ref{app:certificate} includes the complete executable
certificate for the finite calculation.
